\documentclass[xcolor=dvipsnames, compress, 12pt, aspectratio=169, handout]{beamer}
\usepackage{../preamble_slides}

\title{Ideas on Immigration: H-1B Visas and Firms \& Workers}
\author{Amy Kim}
\date{\today}

\begin{document}
\begin{frame}[plain]
    \titlepage
\end{frame}

\begin{frame}
    \frametitle{H-1B Ideas: Overview}
    \begin{itemize}
        \item \textbf{Data:} FOIA H-1B data + LinkedIn Revelio data 
        \begin{itemize}
            \item Have data on full universe of H-1B applicants (lottery winners and losers) starting 2020, from 2009-2019 only data on lottery winners + a proxy for employer H-1B demand (LCAs) used in other papers 
        \end{itemize}
        \item \textbf{Idea One:} Effects of H-1B on Firms \& Workers
        \begin{itemize}
            \item Mixed results on crowd-out and firm patenting/growth from existing literature 
            \item No papers (\textbf{yet}) using new 2020-2024 data (full universe of H-1B applicants)
            \item Can answer old questions (crowd-out of H-1B visas) and new questions (Are H-1B hires higher `quality' than counterfactual hires? If crowd-out of incumbents, where do incumbents go? Effects on return migration?)
        \end{itemize}
        \item \textbf{Idea Two:} Ethnic Sorting of H-1B Workers 
        \begin{itemize}
            \item Effect of winning additional H-1B visa for worker from origin country $c$ on future hiring from $c$
            \item Unclear on how to disentangle mechanisms
        \end{itemize}
    \end{itemize}
\end{frame}

\begin{frame}\label{data}
    \frametitle{Available Data}
    \textbf{2009-2019: FOIA data + Labor Condition Applications (LCAs)}
    \begin{itemize}
        \item Full universe of LCAs (firms must submit approved LCA with H-1B application) \hyperlink{h1bfiling}{\beamergotobutton{Filing Process}}
        \begin{itemize}
            \item Employer information and job information but no employee information 
            \item Also no guarantee that LCA is for H-1B
        \end{itemize}
        \item FOIAed data on H-1B lottery winners: employer name and address, applicant gender, country of origin, \textcolor{gray}{YOB}, \textcolor{gray}{current visa status}, job title, salary, education level, \textcolor{gray}{field of study}, exemption status
        \item Lots of papers use LCAs as proxy for total number of H-1B applications
    \end{itemize}
    \textbf{2020-2024: Bloomberg FOIA data} \hyperlink{h1bdesc}{\beamergotobutton{Descriptives}}
    \begin{itemize}
            \item Full universe of H-1B lottery applications
            \item For losers: only company name and address, applicant gender and country and YOB, preparing `agent' name 
            \item For winners: same data from FOIA above 
    \end{itemize}
\end{frame}

\begin{frame}[allowframebreaks]\label{idea2}
    \frametitle{Idea One: Revisiting Effects of H-1B on Firms and Workers}
    \textbf{Open Questions:} Does H-1B crowd out incumbent workers? Does H-1B improve firm innovation/growth? \hyperlink{h1blit}{\beamergotobutton{Lit Review}}
    \begin{itemize}
        \item Older papers using DiD with H-1B cap changes mostly suggest crowd-in and pos. effects on innovation \hyperlink{h1bpolicy}{\beamergotobutton{H-1B Context}}
        \item Mixed results for recent papers using lottery (different samples, time periods, data)
        \begin{itemize}
            \item Doran, Gelber and Isen (2022 JPE) [DGI] suggest significant crowd-out with no effect on patents; Dimmock, Huang and Weisbenner (2022 Management Science) [DHW] suggest large innovation effects for high-tech firms; Mahajan et al. (2024 WP) [MMSCB] suggest no crowd-out and large pos. eff. on firm growth
        \end{itemize}
    \end{itemize}
    \framebreak 
    \textbf{Potential Contributions:}
    \begin{itemize}
        \item Does not look like anyone has used the new Bloomberg data in a paper? Could use for better estimate of crowd-out.
        \begin{itemize}
            \item Benefits of data: data includes both winners and losers (unlike DHW and MMSCB) and all firms are subject to lottery (unlike DGI)
            \item Seems likely that someone is already working on this (Bloomberg data released last July) and maybe contribution to literature \textit{alone} not very large
        \end{itemize}
        \item Preferred idea: Use LinkedIn data + new \& old FOIA data to explore \textbf{what happens to incumbent workers \& new hires at firms who win the lottery vs. firms who don't}
        \begin{itemize}
            \item If crowd-out via hiring, firms who lose lottery will hire someone else instead. Is the H-1B hire more productive/better (measure via promotions/turnover)?
            \item Suggestive evidence that if H-1B hurts incumbents, it is more-tenured/older skilled workers. If this is true, what happens to them (where do they go? Is this an efficient outcome?)
            \item Return migration: are H-1B hires more likely to remain in US for longer (via green card or H-1B renewal)? E.g. alternative pathway for Canadians (TN visa) has no pathway to green card.
            \item If I can get info on previous H-1B applications: what are effects of winning on first try vs. winning on a later try? (descriptively: what do workers who don't win H-1B immediately do in the intervening years?)
            \item Also concerned that someone else might be working on this? (e.g. Khanna and Morales (2024 WP R\&R at AER) use LinkedIn data to measure effects of H-1B on Indian IT industry)
        \end{itemize}
    \end{itemize} 
\end{frame}


\begin{frame}[allowframebreaks]
    \frametitle{Idea Two: Ethnic Sorting of High-Skill Workers}
    \begin{itemize}
        \item \textbf{Firm-level event study:} effect of number of `excess' H-1B wins for employees from origin country $c$ (relative to expectation based on number of applications for employees from $c$ and overall probability of lottery win) on future H-1B applications/new hires from $c$
        \begin{itemize}
            % \item Controlling for incumbent employee ethnicity country composition, company characteristics (size, industry, etc.), total number of H-1B applications \& applicant country composition
            \item \textbf{Identification:} essentially comparing firm A and firm B, both interested in hiring same number of H-1Bs from country $c$ but firm A randomly wins an extra lottery slot for employee from country $c$
        \end{itemize}
        \framebreak
        \item \textbf{Key Roadblock:} I think this will not be as compelling if I can't decompose mechanisms for ethnic sorting\textemdash `employer-side' (learning, information) vs. `employee-side' (referrals/networks, homophily)
        \begin{itemize}
            \item Could use LinkedIn history to ID employees who plausibly worked/studied together, should have stronger effect if `employee-side' matters (but employer could also be learning about specific workplace/school)
            \item Outcomes within same employer but different workplace locations (shuts down homophily but not networks) 
            \item Heterogeneity by employee quality (look forward to see how quickly employee gets promoted relative to average hire) could ID employer-side?
        \end{itemize}
        \item Other issues:
        \begin{itemize}
            \item Only know country of applicant starting in 2020 (before 2020 only know country of winners, would have to use matching or controls)
            \item Not enough origin country heterogeneity? Issues with fraud?
        \end{itemize}
    \end{itemize}
\end{frame}


\appendix 

\begin{frame}\label{h1bpolicy}
    \frametitle{H-1B Policy Context \& Timeline \hyperlink{data}{\beamerreturnbutton{back}}}
    \begin{itemize}
        \item Cap on H-1B visas 2004-today: 65,000 regular, 20,000 `advanced degree exemption' (ADE)
        \begin{itemize}
            \item Older literature uses increase in cap in 1999 and decrease in 2004 as source of variation
        \end{itemize}
        \item H-1B Lottery 
        \begin{itemize}
            \item H-1B petitions allocated on a first-come, first-served basis starting until cap reached 
            \item Most years (FY2008-9 and FY2014 onwards), cap reached immediately so all new H-1B visas allocated via random lottery 
            \item In other years, random lottery only for applications submitted on last day before cap reached 
        \end{itemize}
    \end{itemize}
\end{frame}

\begin{frame}\label{h1bfiling}
    \frametitle{H-1B Filing Process \hyperlink{data}{\beamerreturnbutton{back}}}
        \begin{itemize}
            \item To hire H-1B worker, firm first required to file LCA with DOL attesting that the job will not replace/harm incumbent workers
            \begin{itemize}
                \item Not tied to a specific employee, must be filed no more than six months before start date
                \item Includes information about job type (occupation code), start and end dates, work location, and salary
            \end{itemize}
            \item Prior to FY2021, firms required to fill out full paperwork for H-1B (including LCA and filing fees of $\sim$\$500 per applicant). Only forms of lottery winners processed, rest thrown out. 
            \item Starting FY2021, firms only required to fill out short form and \$10 filing fee to enter lottery, then only lottery winners required to fill out full paperwork. 
            \begin{itemize}
                \item Bloomberg investigation\textemdash streamlined process caused fraud to skyrocket (multiple submissions of same employee from different companies)
            \end{itemize}
        \end{itemize}
\end{frame}

\begin{frame}\label{h1bdesc}
    \frametitle{H-1B Descriptives (FY2023) \hyperlink{data}{\beamerreturnbutton{back}}}
    \begin{itemize}
        \item 474k applications, 128k lottery wins (likely includes cap-exempt employers) [27\%]
        \item 45k employers submit applications, 23k employers have at least one win [51\%]
        \item \textbf{Rampant Fraud:} 35\% of applications are `duplicates' (same employee named on an application with a different employer) and 3.7k employers submit more applications than their total number of current U.S. employees (representing 38\% of applications)
        \item \textbf{Long Right Tail:} 82\% of employers submit five or fewer applications, top 20 employers (incl. Amazon, Google, Meta, etc.) submit 20\% of applications
        \item \textbf{Limited Origin Country Diversity:} 77\% of applicants from India, 7\% from China, next biggest group Canada (less than 1\%) [similar for lottery winners or excluding `duplicate' applications but more reasonable when restricting to employers with less than 5 applications]
    \end{itemize}    
\end{frame}

\begin{frame}[allowframebreaks]\label{h1blit}
    \frametitle{H-1B Literature Review \hyperlink{idea2}{\beamerreturnbutton{back}}}
    \textbf{Older Estimates} (event study/DiD using changes in H-1B cap from 1990s-early 2000s): 
    \begin{itemize}
        \item City-level estimates: Increases in patenting, native employment and wage \textbf{gains} \textcolor{gray}{[Kerr and Lincoln (2010), Peri, Shih and Sparber (2015a, 2015b)]}
        \item Firm-level estimates: Evidence of crowd-in except maybe older skilled natives, positive eff. on growth \textcolor{gray}{[Kerr, Kerr and Lincoln (2015), Mayda et al. (2023)]}
        \item Structural: Immigration \textbf{decr.} native emp \& wages in CS occs but incr. native welfare and firm profits via innovation \textcolor{gray}{[Bound, Khanna and Morales (2017)]}
    \end{itemize}
    \textbf{Newer Estimates} (using random variation from H-1B lottery): mixed results
    \begin{itemize}
        \item Doran, Gelber and Isen (2022): H-1B crowds out incumbent workers (mostly other immigrants) with no effects on patenting or growth
        \begin{itemize}
            \item Use only FY2006-7 lotteries, so sample is ~3k firms who applied on last day (tend to be larger firms, more educated applicants)
            \item Full universe of applications observed; link to IRS admin data + patents 
        \end{itemize}
        \item Mahajan et al. (2024): H-1B winning firms expand with no effects for native/incumbent workers (except for incumbent high-tenure college-educated, interp. as most substitutable)
        \begin{itemize}
            \item Use FY2008 lottery (all firms subject to lottery)
            \item Do not observe lottery losers (have to impute using LCAs); link to LEHD
        \end{itemize}
        \item Dimmock, Huang and Weisbenner (2022): H-1B winning firms more likely to receive extra VC funding, have more patents/citation
        \begin{itemize}
            \item Use later lotteries with larger sample but restrict to high-tech firms
            \item Also do not observe lottery losers; link to Crunchbase
        \end{itemize}
    \end{itemize}
\end{frame}

% \begin{frame}[allowframebreaks]
%     \label{idea1}
%     \frametitle{Idea One: Effects of Occupational Sorting on Wages}
%     \textbf{Question:} What is the effect $\beta$ of occupational ethnic concentration ($N_{o,d,j}^t$ is the number of workers in occupation $j$ from origin country $o$ in MSA $d$ in year $t$) on wages ($w_{o,d,j}^t$ is the average wage of workers in that cell)?
%     \begin{equation*}
%         w_{o,d,j}^t = \delta_o + \delta_{d,j} + \beta N_{o,d,j}^t + X'_{o,d,j}\gamma + \varepsilon_{o,d,j}
%     \end{equation*}
%     \textbf{Identification:} Obvious endogeneity/reverse causality problem. Can use version of \textcolor{gray}{Burchardi, Chaney and Hassan (2019)} ancestry instrument using destination x occupation cells instead of just destination cells. \hyperlink{fullmodel}{\beamerbutton{Full Model}}
%     \begin{itemize}
%         \item First stage: $N_{o,d,j}^{2010} = \delta_o + \delta_{d,j} + \sum_{t<2010} \alpha_t I_{o, -r(d), -k(j)}^t \frac{I_{-c(o),d,j}^t}{I_{-c(o)}^t} + \eta_{o,d,j}$ where $I^t$ represents immigrant inflow in year $t$, $r(d)$ represents census region of $d$, $k(j)$ represents 1-digit occupation group of $j$, $c(o)$ represents continent of country $o$. 
%         \begin{itemize}
%             \item Idea: instrument for $N$ by allocating past immigrant inflows $I_o^t$ to occupation x destination cells based on overall popularity among all immigrants ($\frac{I_{d,j}^t}{I^t}$ share of all immigrants that go to $d,j$ in $t$). 
%             \item But still a problem if large share of immigrants are from $o$ or large share go to $d,j$. So use `excluded' part of inflow $I_{o, -r(d), -k(j)}$ (inflow from $o$ to MSAs in different census region from $d$ and occupations in different 1-digit group from $j$) and interact with `excluded' part of share $\frac{I_{-c(o),d,j}^t}{I_{-c(o)}^t}$ (share of immigrants from countries in different continents from $o$ going to $d,j$)
%         \end{itemize}
%         \item Exclusion restriction: unobserved factors that affect avg wages in location $d$ and job $j$ for ethnic group $o$ are not systematically related to settlement and job choice of migrants from other continents or flow of immigrants from $o$ to other jobs in other regions.
%         \begin{itemize}
%             \item If you believe BCH2019, I think this is reasonable. \hyperlink{bch2019id}{\beamerbutton{Not sure if I believe BCH2019}}
%         \end{itemize}
%         \item Relevance: historical flows to occupation destination cells actually affect ethnic concentration in occupation destination cells today
%         \begin{itemize}
%             \item This is a bit harder to buy. BCH2019 use decennial census data from 1880-2000, maybe would be more reasonable to use annual counts (e.g. ACS from 2000-present)? 
%         \end{itemize}
%         \item Another issue is power. Especially if using ACS, dividing already small ancestry x destination groups into occupation cells is probably infeasible -- could start with large ancestry groups or large occupation cells and see where that gets me.
%         \item Might require some initial descriptive work to see how persistent occupational niches are in modern data and how they evolve 
%     \end{itemize}
% \end{frame}

% \begin{frame}
%     \frametitle{Idea One Point Five: Aggregate Effects of Occupational Sorting}
%     \textbf{Question:} What is the effect of a more/less occupationally segregated labor market on the immigrant-native wage gap? 
%     \begin{itemize}
%         \item Not clear on best way to approach this question, maybe an extension after previous idea?
%         \item Measure labor market segregation by average `HHI' across all occupations
%         \item Would probably require some sort of structural model 
%         \item Alternatively, could use BCH instruments from Idea One: sum/average instruments across all origin countries and occupations to get predicted local labor market `HHI', use to instrument for actual local labor market `HHI' (but if regression is at the LLM level, too few observations?)
%         % \item NB: different from question `what share of imm-native wage gap can be explained by occupation'? more like -- conditional on occupational mix of ethnicity $e$ nationally and conditional on occupational mix of workers in location $d$, what is effect of redistributing workers such that people are completely segregated by ethnicity between locations vs. redistributing workers such that every location has exact same nationally representative ethnic mix
%     \end{itemize}
% \end{frame}

% \begin{frame}
%     \frametitle{Existing literature on causes of immigrant occ sorting}
%     \begin{itemize}
%             \item Networks: most widely studied, very important historically and today \textcolor{gray}{\footnotesize [Locke 2024, Kerr and Mandorff 2023, Lafortune and Tessada 2016, Patel and Vella 2013]}
%             \item Skills/education: less well-documented, smaller role historically \textcolor{gray}{\footnotesize [Locke 2024, Hanson and Liu 2023]}
%             \item Policy: \textbf{understudied} except specific focus on OPT and H1B changes \textcolor{gray}{\footnotesize [Amuedo-Dorantes and Furtado 2019, Amuedo-Dorantes, Furtado and Xu 2019]}
%     \end{itemize}
% \end{frame}

% \begin{frame}[allowframebreaks]
%     \frametitle{Idea Two: Immigration Policy and Occupational Sorting}
%     \textbf{Question:} How does immigration policy affect occupational choice of immigrants? How does is this effect amplified/mediated via networks? 
%     \begin{itemize}
%         \item Idea would be to do an event study/DiD based on policy changes 
%         \begin{itemize}
%             \item Employment-based immigration (Canada began including employment offers as a criterion in point-based system in ~2010)
%             \begin{itemize}
%                 \item Would be interesting (descriptively) to see change in occupation mix
%                 \item Does employment-based immigration follow the same occupational patterns as pre-existing immigration (i.e. is a company more likely to sponsor you if there are more people from your country working in that occupation/industry already?)
%             \end{itemize}
%             \item Occupation x location-specific immigration pathways (Canada's Provincial Nominee Program) 
%             \begin{itemize}
%                 \item What is distaste for living far from network? (Revealed preference of how many people apply for PNP vs how many people are potentially eligible)
%                 \item Could try to estimate a discrete choice model? 
%             \end{itemize}
%             \item Occupational licensing requirements (e.g. which medical exams carry over to US varies by state and origin country)
%             \begin{itemize}
%                 \item Outcome variable: share of immigrants working in that occupation 
%                 \item Would expect to see the share go up if an occupation becomes less restrictive, not that interesting (unless precise zero)
%                 \item But do networks lower `cost' of licensing? E.g. if it's annoying to get a CPA certification but a lot of people from your country are accountants and can help you through the process, would expect CPA restrictions loosening to matter less for you 
%             \end{itemize}
%         \end{itemize}
%     \item Additional thought: how do children of immigrants make occupational decisions? (Is `occupational following' of parents stronger/weaker than for natives?)
%     \end{itemize}
% \end{frame}

% \appendix 
% \begin{frame}
%     \label{fullmodel}
%     \frametitle{Full Model \hyperlink{idea1}{\beamerbutton{back}}}
%     Modelling $N_{o,d,j}^t$ as:
%     \begin{equation*}
%         N_{o,d,j}^t = a_{o,t} + a_{d,j,t} + b_t N_{o,d,j}^{t-1} + I_o^t(c_t \frac{I^t_{d,j}}{I^t} + d_t \frac{N_{o,d,j}^{t-1}}{N_o^{t-1}})
%     \end{equation*}
%     where 
%     \begin{itemize}
%         \item $I_{o}^t$ is `push' component (total flow of new working-age immigrants from country $o$) 
%         \item $\frac{I^t_{d,j}}{I^t}$ is `pull' component (share of all new working-age immigrants who enter occupation $j$ in MSA $d$)
%         \item $\frac{N_{o,d,j}^{t-1}}{N_o^{t-1}}$ is `recursive' component (share of everyone from $o$ working in $d,j$ last period)
%     \end{itemize}
%     Notes:
%     \begin{itemize}
%         \item Might need to do cells separately by gender?
%         \item Include unemployment as occupation 
%     \end{itemize}
% \end{frame}

% \begin{frame}[allowframebreaks]
%     \label{bch2019id}
%     \frametitle{Issue with BCH2019 Ancestry Instruments? \hyperlink{idea1}{\beamerbutton{back}}}
%     \begin{itemize}
%         \item Suppose some new tech is suddenly discovered in California that both makes it attractive for FDI from China and incentivizes Chinese immigrants to move to California
%         \item Obvious endogeneity problem w naive ols (will observe both high FDI and high ancestry in 2010 but for exog reason), so instrument for ancestry with interaction of $I_o^t\times\frac{I_j^t}{I^t}$. 
%         \item Problem is that Chinese immigrants are large share of all and California arrivals are large share of all. So increase in Chinese immigration to California will inflate both $I_o^t$ and $\frac{I_j^t}{I^t}$ (since more Chinese immigrants are arriving because of California tech and more of all immigrants are going to California because of China). 
%         \item BCH solution is to use $I_{o,-r(d)}^t\times\frac{I_{-c(o),j}^t}{I_{-c(o)}^t}$, interaction of non-West coast Chinese immigrant flow with share of non-Asian immigrants going to California. But if something is pulling Chinese immigrants to California, couldn't that also divert Chinese immigrants from non-West coast places, reducing $I_{o,-r(d)}$? Similarly, if something is pulling Chinese immigrants to California, maybe that will crowd out immigrants from other places, reducing $\frac{I_{-c(o),j}^t}{I_{-c(o)}^t}$?
%         \item So then instrument is correlated with unobserved factors affecting China-California-specific FDI, not excluded?
%     \end{itemize}
% \end{frame}
\end{document}